
\documentclass[12pt]{article}

% DO NOT EDIT THE MARGINS
\usepackage[margin=2.5cm]{geometry}

% Packages
\usepackage{amsfonts, amsmath, amssymb, amsthm}
\usepackage{listings, color, graphicx}
\usepackage{xurl}
\usepackage{cite}
\usepackage[hidelinks]{hyperref}

\begin{document}

% EXECUTIVE SUMMARY ON FIRST PAGE, ON ITS OWN
% 250 WORD LIMIT

\pagestyle{empty}

\begin{center}
{\Large \bfseries Executive Summary}
\end{center}

Your executive summary should go here. It should be no more than 250 words.
Please refer to the Week 3 Worksheet, video, and interactive lecture for more
information on writing an effective executive summary. Note that you do not
have to use the full 250 words. In fact, a shorter executive summary may
be more effective.

\newpage

% MAIN REPORT

\setcounter{page}{1}
\pagestyle{plain}

\section{General guidance}

Your report should go here. It should be no more than 5 sides (please don't
adjust the margins above). You can use sections, subsections, or other ways
of dividing up your report. Don't forget to typeset all of your mathematics
correctly, and reference equations, figures, and tables in the text, e.g.,
with \verb|\ref|.

For example, you can cite references using
\cite{Diffusion,Diffusion1,Diffusion2,NOAA,Numerics,Oenology}.

Please ensure that you've read the Mark Scheme, and also understood the
University's guidance on plagiarism:
\url{https://www.ed.ac.uk/academic-services/students/conduct/academic-misconduct/plagiarism}.

You will receive a `briefing' document, similar to the one you saw for the
Wine Cellar example. Please study it carefully. You may ask Ben Goddard
(\texttt{b.goddard@ed.ac.uk}) for clarification, although the answer may
well be the realistic `I don't know, that's what we've asked you to look at it'.

\newpage

% REFERENCES

\pagestyle{plain}

\addcontentsline{toc}{section}{References}

\bibliographystyle{unsrt}
\bibliography{references}

\newpage

% REFLECTIVE REPORT

\pagestyle{empty}

\begin{center}
{\Large \bfseries Reflective Report}
\end{center}

\textbf{Include up to one page of group reflection on how the project went,
what you learned from the process of undertaking it, and how you think
you can improve in the future.}

Reflection is an important skill. Not only does it help you understand how
to improve, it demonstrates to future employers that you have a desire
and willingness to do so.

This part of the report asks you to reflect, as a group, on the process
(\emph{not} outcomes!) of working on the project. One reason that this is
a group reflection is to ensure that you discuss the process of working
together and collectively learn from it. It is fine for individuals to add
things that specifically affect them (e.g., Ben learned a lot about
debugging from pair programming with Charlotte), but this is most
definitely not the place to moan about people in your team!

You can find lots of good guidance around reflection in the University's
resources:
\url{https://reflection.ed.ac.uk/reflectors-toolkit}.

Note that there is, rightly in some cases, an emphasis on including
references to relevant literature (as you would do in the main body of
the report). \textbf{There is no need to include references here} --
focus on what happened, what you learned, and what you could do better
next time.

Some general tips are:

\begin{itemize}
    \item Be specific. Use concrete examples and relate your experiences
    to the specific project and your roles within it.

    \item Be honest and critical. Acknowledge your strengths and weaknesses
    as a group. Don't be afraid to discuss challenges and mistakes.

    \item Focus on your learning. Demonstrate what you have learned from
    the experience and how it has improved your skills and/or informed
    what you need to work on. Do not include academic things that you
    have learned -- those should be demonstrated in the main report.
\end{itemize}

We would like to see two main areas in the reflection. They can be separate
or intertwined as you see fit.

\noindent \textbf{1: The process of undertaking the project}

Here you should reflect on how the project went, from coming up with your
initial ideas to submitting your final report. As well as being `critical
friends', willing to point out things that could have been improved,
please ensure that you remember the value of things that you found easy
or were very successful.

Make sure that you reflect on how you dealt with challenges, what you
learned from this, and how you can improve your approach. You may want
to include challenges/successes in:

\begin{itemize}
    \item dividing up the work among the group and working together as a team;
    \item encountering dead ends in modelling/numerics;
    \item finding and curating real-world data;
    \item writing, debugging, and validating code;
    \item writing the report, including fitting in (or not) everything you wanted;
    \item ideas you had that you didn't have time to address.
\end{itemize}

This is not an exhaustive list. You should include things that you feel
are important for what you, as a group, experienced and learned during
the project. We are looking for an authentic discussion, not something
shoehorned into some topics that you think you should be addressing.

\medskip

\noindent \textbf{2: Your employability or suitability for future study}

See
\url{https://reflection.ed.ac.uk/reflectors-toolkit/employability}
for more guidance here.

Here you should focus on identifying what skills you have used and
developed throughout the project and build your (collective) sense of
self-awareness. Make sure that you demonstrate how working on the project
has improved your abilities and knowledge. You should also identify any
areas of weakness, or missing skills, that you think you need to develop
as you move towards employment or further study.

If you have no ideas what to include here then have a look at recent job
adverts in areas you'd like to work, and/or go to talk to the Careers
Service:
\url{https://careers.ed.ac.uk/}.

\medskip

\noindent \textbf{A note on language}

See
\url{https://reflection.ed.ac.uk/reflectors-toolkit/producing-reflections/academic-reflections/language}
for more guidance here.

This is still an academic exercise, so be clear and precise. Make sure
that you are as objective as possible -- this is not an opportunity to
complain about members of the group. Be concise; you only have one page.
Make sure that the chronology is clear; a big part of reflection is to
demonstrate learning and improvement.

\end{document}
